% HFBFFT-OMP -- revised manuscript (section-by-section rebuild)
% Class: elsarticle (Computer Physics Communications). Compile with pdflatex.
% Figures in Figures/ ; \graphicspath set below.
\documentclass[preprint,12pt]{elsarticle}
\usepackage{amsmath,amssymb}
\usepackage{graphicx}
\graphicspath{{Figures/}}
\usepackage{booktabs}
\usepackage{siunitx}
\usepackage[T1]{fontenc}
\usepackage{url}
\numberwithin{equation}{section}
\journal{Computer Physics Communications}

\begin{document}

\begin{frontmatter}
\title{HFBFFT-OMP: A Scalable Thread-Parallel 3D Skyrme Hartree--Fock--Bogoliubov Solver}

\author[msu]{Sagar Ashwathanarayana Kikkeri}
\author[msu]{Bailey Knight}
\author[msu]{Kyle Godbey}
\author[msu]{Hasan Metin Aktulga\corref{cor}}
\ead{hma@msu.edu}
\cortext[cor]{Corresponding author}
\address[msu]{Department of Computer Science and Engineering, Michigan State University, East Lansing, Michigan 48824, USA}

\begin{abstract}
HFBFFT solves the Hartree--Fock--Bogoliubov (HFB) equations for atomic nuclei in a three-dimensional coordinate-space representation, the approach best suited to weakly bound and strongly deformed systems. Its shared-memory parallelization distributes the single-particle orbitals across threads. Because the number of orbitals is fixed by the nucleus and does not grow with the core count, this strategy scales poorly on modern many-core nodes. We introduce parallelization strategies that deepen the scaling axis from the fixed orbital count to the grid size, so that the available concurrency grows with the numerical resolution. As a result, the solver keeps improving from the 32--48 cores at which the baseline saturates to 96--128 cores of a 192-core node, while reproducing the converged binding energies exactly.
\end{abstract}

\begin{keyword}
Nuclear density functional theory \sep Hartree--Fock--Bogoliubov method \sep Skyrme energy density functional \sep Parallel algorithms \sep Shared-memory parallelization \sep OpenMP \sep Cache locality
\end{keyword}
\end{frontmatter}

% ======================================================================
\section{Introduction}\label{sec:intro}

Nuclei far from the line of stability are a central subject of nuclear structure research~\cite{erler2012}. Their weak binding couples bound states to the particle continuum, which requires the full Hartree--Fock--Bogoliubov (HFB) treatment of pairing rather than the simpler BCS approximation~\cite{dobaczewski1984}, and their diffuse densities and often exotic shapes require large coordinate boxes with no imposed spatial symmetry~\cite{bender}. These requirements are met by solving the HFB equations on a three-dimensional Cartesian grid.

HFBFFT~\cite{hfbfft} is such a solver. It works in the canonical basis on a 3D grid, evaluates derivatives with fast Fourier transforms, and imposes no spatial symmetry; it is built on the TDHF code Sky3D~\cite{sky3d}, from which it inherits both its numerical methods and its parallelization. Basis-expansion solvers such as HFBTHO~\cite{hfbtho} and HFODD~\cite{hfodd} are more efficient for well-bound nuclei but lose accuracy for weakly bound systems, where large configuration spaces are needed to represent the asymptotic densities. The coordinate-space approach avoids this at a substantial computational cost: each self-consistent iteration operates on every canonical state over the entire grid, and convergence requires thousands of iterations.

\paragraph{The parallelism deficit}
One thing has changed since the parallelization of this solver family was designed. Cores per socket have grown roughly sixfold --- from 16 on the reference node of the earlier distributed-memory Sky3D study~\cite{sky3dmpi} to 96 on a current AMD EPYC 9654~\cite{amdgenoa}. The number of orbitals, by contrast, is fixed by the nucleus and the pairing window, not by the machine. The scheme does also rely on the internal multithreading of the BLAS and FFT libraries, but at the matrix and transform shapes that arise here this contributes little (Section~\ref{sec:perphase}). The solver therefore no longer makes effective use of a full node: the phases that dominate the run time stop improving well before the available cores are exhausted.

\paragraph{Contributions}
This work removes that limitation by exposing, in each phase, a form of concurrency whose extent grows with the numerical resolution rather than with the particle number. Our contributions are:
\begin{itemize}
\item \textbf{Constraint step [\texttt{diagstep}].} We reformulate the two matrix products of the canonical-basis constraint step --- the most expensive phase --- as tall-and-skinny matrix multiplications blocked along their long dimension~\cite{ballard2011}: the Gram product blocks the contraction, and the reconstruction blocks the output. This is the single largest source of the overall speedup.
\item \textbf{Mean-field application [\texttt{grstep}].} We decompose the spectral derivative of each orbital into spatial slabs, mapped to a cooperative team of threads pinned within one shared cache domain.
\item \textbf{Density accumulation [\texttt{add\_density}].} We replace the runtime reduction over the shared density fields, which suffers from false sharing and reduction traffic that grow with the thread count, with a two-stage, output-indexed reduction that removes both.
\item \textbf{Isospin concurrency [\texttt{diagstep} + \texttt{pairing}].} We run the neutron and proton computations concurrently in the constraint and pairing phases, an additional axis of extent two.
\end{itemize}

Section~\ref{sec:phys} develops the physics, the structure of an iteration, and the Fortran-level computation of each phase we touch. Section~\ref{sec:base} explains why the baseline runs out of parallelism, and frames Section~\ref{sec:alg}, which presents the new algorithms. Section~\ref{sec:eval} evaluates them, and Section~\ref{sec:concl} concludes.

% ======================================================================
\section{Physics and computational structure}\label{sec:phys}

This section sets out the physics HFBFFT solves (Sections~\ref{sec:edf}--\ref{sec:canon}) and the numerical representation, structure, and cost of one self-consistent iteration (Section~\ref{sec:iter}). It is written to be read by both computational and nuclear-physics audiences. A complete description of the method is given in Ref.~\cite{hfbfft}.

\subsection{The Skyrme energy density functional}\label{sec:edf}
HFBFFT describes a nucleus by a set of single-particle wavefunctions. We call each one a \emph{state}, or interchangeably an \emph{orbital}: a complex two-component spinor $\psi_\alpha(\mathbf r,s)$ defined over the entire grid and indexed by spin $s$, with $\Omega_q$ of them for each isospin $q\in\{n,p\}$ (neutrons, protons). From the orbitals one builds \emph{fields} --- quantities defined at every grid point. The primary fields are a small set of local densities~\cite{vautherin1972,bender}:
\begin{align*}
\rho_q(\mathbf r) &= \sum_{\alpha\in q}\sum_s v_\alpha^2\,|\psi_\alpha(\mathbf r,s)|^2 && \text{(particle density),}\\
\tau_q(\mathbf r) &= \sum_{\alpha\in q}\sum_s v_\alpha^2\,|\nabla\psi_\alpha(\mathbf r,s)|^2 && \text{(kinetic density),}\\
\mathbf J_q(\mathbf r) &= -i\sum_{\alpha\in q}\sum_{ss'} v_\alpha^2\,\psi_\alpha^*(\mathbf r,s)\,\nabla\times\boldsymbol\sigma_{ss'}\psi_\alpha(\mathbf r,s') && \text{(spin-orbit density),}
\end{align*}
each an occupation-weighted sum over all orbitals, with amplitudes $v_\alpha$ (the occupations, defined in Section~\ref{sec:canon}). The total energy is a functional of these density fields,
\[
E_{\mathrm{tot}} = E_{\mathrm{Skyrme}}[\rho,\tau,\mathbf J] + E_{\mathrm{pair}}[\rho,\xi] + E_{\mathrm{Coul}}[\rho_p],
\]
combining the Skyrme functional in the particle--hole channel (the SLy4 parametrization here~\cite{chabanat1998}), a Coulomb term acting on the proton density $\rho_p$, and a density-dependent contact interaction in the particle--particle channel, which acts through the pairing density $\xi_q$ built from the products $u_\alpha v_\alpha$~\cite{dobaczewski1984}. Differentiating this energy with respect to the densities yields a second kind of field --- the \emph{potentials} that act back on the orbitals --- which drive the minimization in Section~\ref{sec:canon}.

One structural point carries forward: every density field is a sum over all orbitals of a pointwise quantity on the grid. This is the origin of the density-accumulation phase, and the reason that phase is a reduction.

\subsection{HFB in the canonical basis}\label{sec:canon}
The states and occupations of Section~\ref{sec:edf} are found by \emph{minimizing} $E_{\mathrm{tot}}$. The two potentials named there are its functional derivatives --- the mean field $\hat h=\delta E/\delta\rho$ and the pairing field $\hat{\tilde h}=\delta E_{\mathrm{pair}}/\delta\xi$. HFBFFT performs the minimization in the \emph{canonical basis}, the representation in which the one-body density matrix is diagonal and the HFB state reduces to independent BCS-like pairs $(\alpha,\bar\alpha)$ with occupations $v_\alpha^2$ and complementary amplitudes $u_\alpha^2=1-v_\alpha^2$~\cite{ringschuck}. Varying the energy with respect to $\psi_\alpha$ and $v_\alpha$ then gives the mean-field equations
\begin{equation}\label{eq:mf}
\hat H_\alpha\psi_\alpha=\sum_\beta\psi_\beta\lambda_{\beta\alpha},\qquad
\hat H_\alpha = v_\alpha^2\,\hat h + u_\alpha v_\alpha\,\hat{\tilde h},
\end{equation}
where $\lambda$ is the matrix of Lagrange multipliers enforcing orthonormality of the states,
\begin{equation}\label{eq:lambda}
\lambda_{\beta\alpha}=\tfrac12\langle\psi_\beta|\hat H_\alpha+\hat H_\beta|\psi_\alpha\rangle.
\end{equation}
The occupations are fixed by the gap equation, which has the closed form
\begin{equation}\label{eq:gap}
v_\alpha^2=\tfrac12-\tfrac12\,\frac{h_{\alpha\alpha}-\epsilon_{F,q}}{\sqrt{(h_{\alpha\alpha}-\epsilon_{F,q})^2+\Delta_{\alpha\alpha}^2}},
\end{equation}
with $h_{\alpha\alpha}$ and $\Delta_{\alpha\alpha}$ the diagonal matrix elements of $\hat h$ and $\hat{\tilde h}$, and $\epsilon_{F,q}$ the Fermi energy set by the particle number.

Convergence is monitored through the residual of the symmetry condition $\lambda^-_{\beta\alpha}=0$,
\begin{equation}\label{eq:ds}
\Delta S=\sqrt{\tfrac12\sum_q\frac1{\Omega_q^2}\sum_{\alpha,\beta\in q}\bigl|\lambda^-_{\beta\alpha}\bigr|^2},\qquad
\lambda^-_{\beta\alpha}=\tfrac12\langle\psi_\beta|\hat H_\alpha-\hat H_\beta|\psi_\alpha\rangle,
\end{equation}
which vanishes at self-consistency. Section~\ref{sec:correct} uses $\Delta S$ to verify that the reformulated constraint step reproduces the original iteration.

\subsection{Numerical representation and structure of an iteration}\label{sec:iter}
The equations of Section~\ref{sec:canon} are coupled --- the fields depend on the densities, which depend on the states --- and are solved by a self-consistent fixed-point iteration. We first fix the numerical representation, then walk one iteration, giving for each phase the computation it performs and its dominant cost.

\paragraph{Representation}
All fields live on an $n_x\times n_y\times n_z$ Cartesian grid with no imposed symmetry. A canonical state is a complex spinor stored in Fortran as \texttt{psi(nx,ny,nz,2)}, so the $x$ index is contiguous and the spin index is outermost. Flattened, one state is a vector of length $M=2\,n_xn_yn_z$, and the $\Omega_q$ states of one isospin form the \emph{tall-and-skinny} matrix $\psi_q$ of shape $M\times\Omega_q$ with $M\gg\Omega_q$. Derivatives are spectral: along a given direction, a forward FFT, a multiplication of each mode by $(ik)^m$, and an inverse FFT. For the first derivative along $x$ the factor is $ik_x$, with $k_x=2\pi\nu/(n_x\,dx)$ the wavenumber of mode $\nu$; along $x$ the state is a batch of $2\,n_yn_z$ independent lines of length $n_x$, one per $(y,z,\text{spin})$.

\paragraph{One iteration}
Damped gradient iteration~\cite{blum} repeats the following steps until $\Delta S$ (Eq.~\eqref{eq:ds}) converges. The four phases we optimize are set in bold, with the timer names used in the rest of the paper to refer to each phase denoted in square brackets. Figure~\ref{fig:flow} shows the resulting flow.
\begin{enumerate}
\item \textbf{Density accumulation [\texttt{add\_density}]} --- accumulate $\rho,\tau,\mathbf J,\xi$ as the orbital sums of Section~\ref{sec:edf}, then build the fields $\hat h,\hat{\tilde h}$. This amounts to a reduction into the shared grid fields through a pass over all $\Omega_q$ states.
\item \textbf{Mean-field application [\texttt{grstep}]} --- apply $\hat h,\hat{\tilde h}$ to every state and take the damped gradient step. The cost is dominated by the spectral derivatives ($\nabla$ and $\nabla^2$) of each state, i.e.\ $O(\Omega_q)$ batched FFTs, plus pointwise multiplications by the fields.
\item \textbf{Constraint step [\texttt{diagstep}]} --- form the $\Omega_q\times\Omega_q$ matrices $\psi_q^H\psi_q$ and $\psi_q^H(\hat h\psi_q)$, diagonalize, and apply the resulting transform back to the states ($\psi_q U$). The state dependence of $\hat H_\alpha$ in Eq.~\eqref{eq:mf} makes the multiplier matrix $\lambda$ (Eq.~\eqref{eq:lambda}) \emph{full}, so the entire $\Omega_q\times\Omega_q$ matrix must be assembled and applied every iteration --- unlike Hartree--Fock or HF+BCS, including Sky3D, which need only its diagonal. The two length-$M$ contractions ($M\times\Omega_q\to\Omega_q\times\Omega_q$) and the reconstruction ($M\times\Omega_q$ by $\Omega_q\times\Omega_q$) essentially account for all the arithmetic in this phase with the cost of  $\Omega_q\times\Omega_q$ eigensolve being negligible. Overall, this is the most computationally expensive phase of an iteration.
\item \textbf{Pairing [\texttt{pairing}]} --- solve the gap equation~\eqref{eq:gap} for the occupations. For open-shell nuclei this runs a configuration-space sub-iteration: many small dense matrix sweeps (a gradient update, a QR orthonormalization, a change of basis) applied in strict sequence --- a recurrence --- followed by one large recombination of the states. It is interleaved with the 3D steps to reduce their number.
\item \textbf{Convergence test} on $\Delta S$.
\end{enumerate}

Two parameters, a soft cutoff on the pairing-active space and pairing annealing, change the size of the problem ($\Omega_q$ rises above the number of particles present for a given nucleus), but the overall computational structure remains the same for different values of these parameters.

Taken together, the constraint step (\texttt{diagstep}) dominates an iteration, followed by the mean-field application (\texttt{grstep}). For open-shell nuclei, the pairing sub-iteration can become expensive. Typically, density accumulation (\texttt{add\_density}) accounts for a small fraction of the overall runtime.  Section~\ref{sec:eval} gives a detailed breakdown of the measured timings for sequential as well as parallel executions under different scenarios.

\begin{figure}[tbp]\centering
  \includegraphics[width=0.7\textwidth]{Figure1}
  \caption{One HFBFFT iteration, annotated with the axis of parallelism each phase exposes after the redesign of Section~\ref{sec:alg}. Shading marks the scheme: orbital-parallel with an output-indexed reduction, cooperative spatial teams, blocked matrix kernels with isospin teams, recurrence-bound with isospin teams, and serial or redundant.}
  \label{fig:flow}
\end{figure}

% ======================================================================
\section{Why the baseline parallelization runs out of work}\label{sec:base}

The baseline is the OpenMP version of HFBFFT as published~\cite{hfbfft}. Its structure is uniform across phases: a single \texttt{OMP DO} over the orbitals, one thread per orbital, a parallel region opened and closed per kernel, a shared-array reduction for the densities, and the two isospins processed one after the other. Where a phase's heavy arithmetic is a matrix product or a transform, it is handed to a single BLAS or FFTW call at default threading. The only axis the code itself exposes is the orbital index, whose extent $\Omega_q$ is fixed by the nucleus and does not grow with the core count.

The evaluation platforms differ chiefly in one property --- their last-level-cache organization, characterized by the number of cores sharing a cache domain, $D$, and that domain's private capacity, $L$. This is the structure the second failure mode below is sensitive to.

A phase fails to use a growing core count in one of two ways.
\begin{itemize}
\item \textbf{Exhaustion} --- the phase does not improve with added cores because no axis is effectively distributed across them. This has two causes. In the constraint step (\texttt{diagstep}) the work exists but is hidden inside single black-box BLAS calls that partition only the product's output: for the Gram product that output is tiny ($\Omega_q\times\Omega_q$), and for the reconstruction it is large but threaded without cache or NUMA locality, so the phase saturates by about 8 cores. This is \emph{curable} --- Section~\ref{sec:constraint} takes control of the blocking. In the pairing sub-iteration, by contrast, no axis exists at all: the inner sweeps are a genuine sequential recurrence, so the exhaustion is \emph{incurable}.
\item \textbf{Inversion} --- the phase does expose a growing axis and improves up to a threshold, beyond which more cores make it slower because they saturate a shared resource. The mean-field application (\texttt{grstep}) inverts when the working sets of co-resident orbitals overflow a cache domain (Section~\ref{sec:mf}); the density accumulation (\texttt{add\_density}) inverts when threads contend on the shared output through false sharing and reduction traffic (Section~\ref{sec:dens}). Both set in within a single socket; the NUMA boundary only compounds them.
\end{itemize}

Table~\ref{tab:dop} summarizes the degree of parallelism (DoP) each phase exposes, before and after the reformulations of Section~\ref{sec:alg}. The redesign's axes are independent of the orbital index and grow with the numerical resolution --- $M$ and $n_z$ increase with grid size and box length while $\Omega_q$ is fixed --- so they multiply with the orbital axis and lengthen as a calculation is refined.

\begin{table}[tbp]\centering\small
\caption{Degrees of parallelism exposed by each phase, before and after the redesign. $b_K,b_R$ are the contraction and output blocking factors of Section~\ref{sec:constraint}. $\dagger$~The baseline \texttt{diagstep} parallelism is nominal only: it is provided by BLAS's internal multithreading, which black-boxes the work partitioning, and at these tall-skinny shapes it saturates by ${\approx}8$ threads --- the measured \texttt{diagstep} plateau. FFTW similarly threads \texttt{grstep} internally.}
\label{tab:dop}
\begin{tabular}{lclc}\toprule
Phase & Baseline axis (DoP) & Axes after redesign & DoP\\\midrule
\texttt{grstep} & orbitals ($\Omega_q$)$\dagger$ & orbitals $\times$ spatial slabs & $\Omega_q\cdot 2n_z$\\
\texttt{diagstep} Gram & $O(\Omega_q^2)$$\dagger$ & isospin $\times$ contraction blocks & $2\lceil M/b_K\rceil$\\
\texttt{diagstep} reconstruction & $O(M\Omega_q)$$\dagger$ & isospin $\times$ output blocks & $2\lceil M/b_R\rceil$\\
\texttt{pairing} recurrence & none ($1$) & isospin & $2$\\
\texttt{pairing} recombination & orbitals ($\Omega_q$) & isospin $\times$ output blocks & $2\lceil M/b_R\rceil$\\
\texttt{add\_density} accumulate & orbitals ($\Omega_q$) & orbitals (unchanged) & $\Omega_q$\\
\texttt{add\_density} merge & tree reduction & output planes & $n_z$\\\bottomrule
\end{tabular}
\end{table}

% ======================================================================
\section{Parallel algorithms}\label{sec:alg}

\subsection{Constraint step: blocking along the long dimension}\label{sec:constraint}
The constraint step assembles and applies the multiplier matrix $\lambda$ through two matrix products with the tall-skinny state block $\psi_q$ ($M\times\Omega_q$, $M\gg\Omega_q$): the \emph{Gram product} $\psi_q^H(\hat h\psi_q)$, which contracts the long dimension down to a small $\Omega_q\times\Omega_q$ matrix, and the \emph{reconstruction} $\psi_q U$, which applies the resulting $\Omega_q\times\Omega_q$ transform back to produce a new $M\times\Omega_q$ block. As Section~\ref{sec:base} explained, the baseline computes each as a single BLAS call that distributes only the product's output~\cite{blis-mt}, and so does not scale (Table~\ref{tab:dop}).

In both products the arithmetic lies along the long $M$ dimension, and the reformulation blocks that dimension explicitly. Because $M$ is the \emph{contraction} of the Gram product but the \emph{output} of the reconstruction, the same idea takes two forms.

\paragraph{Gram product --- contraction blocking}
We divide the $M$ contraction rows into blocks of $b_K$. Each thread takes a block, forms the partial $\Omega_q\times\Omega_q$ product of its slice with a single small ZGEMM into a private buffer, and the private buffers are then summed into the shared result in one critical region. Only that $\Omega_q^2$ combine is serialized; everything else runs in parallel. This exposes $\lceil M/b_K\rceil$ independent blocks per isospin, where the black-box call exposed almost none.

\paragraph{Reconstruction --- output blocking}
We divide the $M$ output rows into blocks of $b_R$ and give each thread a block, which it fills with a single small ZGEMM writing its own disjoint rows. No reduction is needed, and because each thread writes a contiguous range of rows its data stays within one cache and NUMA domain, removing the memory-locality penalty of the baseline. This exposes $\lceil M/b_R\rceil$ blocks per isospin.

The $\Omega_q\times\Omega_q$ eigen-decomposition between the two products is left as a single LAPACK call; at these $\Omega_q$ it is negligible. Following the distributed-memory Sky3D work~\cite{sky3dmpi}, we form the product of the orthonormalization and diagonalization matrices first and apply the single combined transform to the states, which halves the number of passes over the full $M\times\Omega_q$ block. The blocking factors are fixed at $b_K=512$ and $b_R=1024$ throughout (Section~\ref{sec:setup}).

\subsection{Isospin concurrency}\label{sec:iso}
Neutrons and protons couple only through the shared mean field. Once $\hat h$ and $\hat{\tilde h}$ have been built at the start of an iteration, the constraint step and the pairing phase for the two isospins are independent of each other until the next density accumulation. We exploit this by running the two isospins concurrently: the constraint step (\texttt{diagstep}) runs them on two separate teams, each using half the available cores, and the pairing phase (\texttt{pairing}) runs a two-way loop over isospin. The axis is confined to these two phases --- \texttt{grstep} and \texttt{add\_density} either couple the isospins through the shared field or are already orbital-parallel. The same neutron/proton split is used in the distributed-memory Sky3D solver~\cite{sky3dmpi}, there apportioned across MPI ranks in proportion to the particle numbers.

The two phases use different splitting mechanisms, because of how libgomp handles nesting. A top-level OpenMP region reuses a persistent \emph{warm} pool --- the team is created once and thereafter only woken --- whereas a nested, second-level region gets \emph{cold} threads, created afresh on every entry. In the constraint step the matrix kernels of Section~\ref{sec:constraint} are entered about nine times per isospin each iteration; splitting the isospins with OpenMP would nest those kernels, so at 192 cores every entry spawns roughly 95 cold workers --- of order 1{,}700 thread creations per iteration --- which roughly doubled the constraint-step kernel time. We split the isospins with POSIX threads instead: two pthreads per diagonalization, inside which each isospin runs its kernels as top-level, warm-pool regions. That is two thread creations per diagonalization rather than thousands.

The pairing phase keeps a plain OpenMP two-way loop. Its only multithreaded kernel, the tall-skinny recombination, runs once at the end of an otherwise serial sub-iteration, so nesting would create cold threads just twice per iteration --- too rarely to matter, and cheaper than managing pthreads.

\subsection{Mean-field application: slab decomposition}\label{sec:mf}
Applying $\hat h$ and $\hat{\tilde h}$ to a state (\texttt{grstep}) is dominated by its spectral derivatives. A derivative along any grid direction is three dependency-ordered stages: a forward FFT along that direction, a multiplication of each Fourier mode by $i$ times its wavenumber along that direction, and an inverse FFT, each stage consuming the full output of the previous. The baseline evaluates each derivative as a single batched transform over all of the state's lines in that direction (Section~\ref{sec:iter}), so the derivative is one indivisible unit of work and the only concurrency is across orbitals --- the axis that does not grow (Section~\ref{sec:base}).

Splitting a derivative stage by stage across a team would not help, because each stage needs the whole array from the previous one, which would force the team to synchronize three times per derivative. We regroup the work instead. Taking the $x$-derivative as the example, its $2n_yn_z$ length-$n_x$ transforms are partitioned into $2n_z$ \emph{slabs}, one per (spin, plane); a parallel loop assigns each slab to a thread, and each thread carries its slab through the forward transform, the modal multiplication, and the inverse transform in sequence. The dependency is then satisfied privately within the slab, and the team synchronizes only once, at the end of the loop; the $y$ and $z$ derivatives are decomposed the same way, by their own planes. Each orbital's derivative is thereby split into $2n_z$ independent work items, where the baseline exposed none.

Two properties follow from partitioning by output. Each thread writes a disjoint output slab, so no output memory is shared and nothing races; and the transform scratch is declared inside the loop body, giving each thread a private buffer rather than a shared one. How many threads cooperate on one orbital --- the team width $w$ --- is set by the cache and fixed in Section~\ref{sec:mapping}.

\subsection{Density accumulation: output-indexed reduction}\label{sec:dens}
The orbital axis is already exposed here --- the densities are sums over orbitals (Section~\ref{sec:edf}). What does not scale is \emph{combining} the per-thread contributions. The phase accumulates six shared real grid fields --- the particle, kinetic and pairing densities, the current, and the spin and spin-orbit densities --- over all orbitals, and in the baseline every thread adds into the shared fields through an OpenMP reduction. Two costs then grow with the thread count: false sharing, because adjacent grid cells written by different threads fall on shared cache lines (Section~\ref{sec:base}), and the reduction itself, since the runtime gives each thread a private copy of every full field and combines them at a barrier.

We replace this with a two-stage, output-owned scheme. In the first stage each thread accumulates its orbitals into its own preallocated copy of the partial fields --- ordinary privatization. In the second stage a parallel loop runs over the output planes: each thread owns one $z$-plane of every field, zeroes it, and sums that plane across all threads' partials. Because the loop is indexed by the output and distinct planes lie on disjoint cache lines, no reduction clause and no atomics are needed, and nothing is write-shared. This exposes $n_z$ independent work items, independent of the thread count. The cost is memory --- one private copy of the fields per thread --- which is why this phase's thread count is capped (Section~\ref{sec:setup}).

\subsection{Cache-aware mapping: team width}\label{sec:mapping}
Modern server CPUs do not present their cores as one flat pool behind a single cache. The cores are grouped into clusters, and each cluster has its own last-level (L3) cache, local to it and shared among its cores. We call such a cluster a \emph{cache domain}: $D$ cores backed by a local cache of capacity $L$, with several domains to a socket. Data used by a core is fastest when it sits in that core's own domain cache.

This structure is why the mean-field application (\texttt{grstep}) has several threads cooperate on a single orbital rather than simply running one orbital per core. The slab decomposition of Section~\ref{sec:mf} makes an orbital's derivative divisible, so a team of threads can share it --- nested parallelism. It is tempting to expose as many independent orbitals as there are cores, but an orbital's working set is large, and the orbitals worked on together within one domain must fit that domain's cache. Cooperating $w$ threads on each orbital reduces how many orbitals are resident in a domain at once, keeping them within $L$.

This fixes the team width. The number of orbitals a domain can hold at once is its capacity divided by one orbital's \emph{hot} footprint --- the data actually live at each step,
\begin{equation}\label{eq:cap}
n \;=\; L/F_{\text{hot}},
\end{equation}
and the team width follows from sharing the domain's $D$ cores among them,
\begin{equation}\label{eq:width}
w \;=\; D/n \;=\; D\,F_{\text{hot}}/L.
\end{equation}
The hot footprint is only a fraction of an orbital's full working set; it, and the resulting optimum width, are quantified in Section~\ref{sec:capacity}. On machines with small private per-domain caches the bound is active and fixes a definite width; on machines whose last-level cache is shared across a socket it does not bind ($L$ is large), and $w$ is set instead by transform-batch efficiency and NUMA placement, so the algorithm is architecture-independent even where the width is not.

Placement then follows the same two-level structure: cooperating threads are packed within one cache domain so a team shares a single local cache, while separate teams are spread across domains; OpenMP nesting is held to two active levels, the orbital teams above and the slab threads within.

% ======================================================================
\section{Performance evaluation}\label{sec:eval}

\subsection{Experimental setup}\label{sec:setup}
\paragraph{Machines}
The primary machine is an exclusively allocated dual-socket AMD EPYC 9654 node (Zen 4, Genoa): $2\times96=192$ physical cores in eight-core core-complexes, each with a private 32 MB L3, two NUMA domains, DDR5-4800 memory. Portability is assessed on the Zen 2 and Skylake nodes of Table~\ref{tab:machines}. Threads are bound at core granularity with \texttt{OMP\_PLACES=cores} and a two-level \texttt{OMP\_PROC\_BIND=spread,close}, and nesting is capped at \texttt{OMP\_MAX\_ACTIVE\_LEVELS=2}.

\begin{table}[tbp]\centering
\caption{Evaluation platforms. The first column is the reference node of the earlier distributed-memory Sky3D study~\cite{sky3dmpi,intel_haswell}. $D$ is the number of cores sharing one cache domain; $L$ is that domain's private capacity.}
\label{tab:machines}
\begin{tabular}{lcccc}\toprule
 & Haswell (ref.) & Skylake & Zen 2 & Zen 4\\
 & E5-2698 v3 & (intel18) & (amd20) & EPYC 9654\\\midrule
Cores / socket & 16 & 20 & 64 & 96\\
Cores / node & 32 & 40 & 128 & 192\\
Cache domain $D$ & socket-wide & socket-wide & 4 cores & 8 cores\\
Domain cache $L$ & 40 MB shared & shared, non-incl. & 16 MB & 32 MB\\
Path between domains & --- & mesh & off-die (I/O die) & Infinity Fabric\\
NUMA domains / node & 2 & 2 & 2 & 2\\\bottomrule
\end{tabular}
\end{table}

\paragraph{Build}
The code is compiled with GCC 12.3.0 at \texttt{-O3 -march=native -ffast-math -funroll-loops -fopenmp}, against FFTW~3~\cite{fftw} and OpenBLAS~0.3.23~\cite{openblas}.

\paragraph{Problem}
All runs use a $48^3$ Cartesian grid at 0.8~fm spacing (a 38.4~fm box), the SLy4 Skyrme functional~\cite{chabanat1998} in the particle--hole channel, and a mixed density-dependent $\delta$ interaction in the pairing channel. The test set separates the two regimes: $^{208}$Pb and $^{132}$Sn are doubly magic and unpaired, isolating the three phases common to every calculation, while $^{120}$Sn, $^{102}$Zr and $^{110}$Zr are open-shell and paired, adding the sub-iteration. The active spaces summed over both isospins are $\Omega_n+\Omega_p=208$ for $^{208}$Pb and $302$ for $^{102}$Zr (the per-isospin $\Omega_q$ entering the kernels is roughly half). For the paired cases the sub-iteration switches on at outer iteration 40 and performs 100 inner sweeps per outer iteration.

\paragraph{Parameters and threading}
These choices instantiate the symbols used above: $M=2\cdot48^3=221{,}184$, $n_z=48$ (so the slab axis has extent $2n_z=96$), and the contraction axis $\lceil M/b_K\rceil=432$ at the blocking factor $b_K=512$; the output blocking factor is $b_R=1024$ and the team width is $w=4$. Thread counts are set per phase: the mean-field application and the constraint step use the full node, the pairing recombination 96, and the density accumulation is capped at 64 --- beyond that its merge stage gains nothing while the per-thread partial fields cost memory, one private copy of the six density fields per thread. The FFT library runs single-threaded within a team.

\paragraph{Protocol}
Reported times are steady-state per-iteration wall-clock times: at each core count we run 150 iterations with per-phase timers and average over the post-warmup window, discarding the first 15. Physics correctness is verified separately by running the highest core count to the input deck's convergence criterion.

\subsection{Per-phase scaling}\label{sec:perphase}
Figures~\ref{fig:grstep}--\ref{fig:add} show the strong scaling of each phase from 1 to 192 cores, baseline and optimized overlaid, one panel per test case. Table~\ref{tab:phase} gives the fastest time and where it is reached; Table~\ref{tab:scaling} reports the more telling quantities --- the core count at which each phase stops improving, and how much it degrades beyond.

\begin{figure}[tbp]\centering
  \includegraphics[width=\textwidth]{figure2}
  \caption{Mean-field application (\texttt{grstep}): per-iteration time versus core count, baseline vs.\ optimized, one panel per test case. The open circle marks the fastest configuration.}
  \label{fig:grstep}
\end{figure}

\begin{figure}[tbp]\centering
  \includegraphics[width=\textwidth]{figure3}
  \caption{Constraint step (\texttt{diagstep}): per-iteration time versus core count, baseline vs.\ optimized, one panel per test case. The baseline flattens at about 8 cores and gains only ${\approx}1.3\times$ to 192, while the optimized version scales across the full node; the open circle marks the fastest configuration.}
  \label{fig:diag}
\end{figure}

The two failure modes of Section~\ref{sec:base} are visible. The constraint step (\texttt{diagstep}, Fig.~\ref{fig:diag}) flattens rather than inverts: over the 8$\to$192-core range, a 24$\times$ increase in resources, the baseline gains only about 1.3$\times$, the signature of work the library cannot distribute; after the reformulation it scales roughly 34$\times$ from 1 to 192 cores. The mean-field application (\texttt{grstep}) and the density accumulation (\texttt{add\_density}) invert. \texttt{add\_density} is the most extreme --- the baseline is fastest at 16 cores and nearly five times slower at 192, as false sharing and reduction traffic overwhelm the useful work --- while the redesign improves to 48 cores and is then flat. The pairing phase (Fig.~\ref{fig:pairing}) is limited by its serial recurrence and improves little in either build.

\begin{figure}[tbp]\centering
  \includegraphics[width=0.55\textwidth]{figure4}
  \caption{Pairing phase (\texttt{pairing}), $^{102}$Zr, baseline vs.\ optimized. The sequential recurrence limits both; the optimized build shifts the optimum from 48 to 64 cores and reduces the degradation beyond it.}
  \label{fig:pairing}
\end{figure}

\begin{figure}[tbp]\centering
  \includegraphics[width=\textwidth]{figure5}
  \caption{Density accumulation (\texttt{add\_density}): per-iteration time versus core count, baseline vs.\ optimized, one panel per test case. The baseline is fastest at 16--32 cores and much slower at 192, driven by false sharing and reduction traffic; the optimized version improves to 48 cores and is then flat. The open circle marks the fastest configuration.}
  \label{fig:add}
\end{figure}

The constraint-step gain separates into two parts. At a fixed 32 cores the $^{208}$Pb time falls from about 0.83~s to 0.31~s, a factor of ${\sim}2.7$ from the reformulated kernels together with isospin concurrency; the redesigned phase then continues to scale from 0.31~s at 32 cores to 0.107~s at 192, a further ${\sim}2.9\times$ from the newly exposed parallelism. (We do not separate the kernel reformulation from isospin concurrency within the first factor.)

Table~\ref{tab:budget} gives the iteration budget at a single core count. In the baseline the constraint step is by a wide margin the largest phase, and the largest single source of savings --- about two-thirds of the end-to-end gain in both nuclei. After the redesign the ordering changes: the mean-field application becomes the largest remaining phase, because the constraint step improved so much more.

\begin{figure}[tbp]\centering
  \includegraphics[width=\textwidth]{figure6}
  \caption{End-to-end per-iteration time, baseline vs.\ optimized, one panel per test case. The baseline turns upward at 32--48 cores, well inside a single 96-core socket; the optimized version descends monotonically, crosses the socket boundary without penalty, and is flat from 96 to 192 cores. The curves also cross below 16 cores, where the fixed overhead of the added axes is not yet amortized.}
  \label{fig:e2e}
\end{figure}

\begin{table}[tbp]\centering
\caption{Fastest per-phase time and the core count at which it is reached, on the EPYC 9654 node. Speedups are from full-precision times.}
\label{tab:phase}
\begin{tabular}{llccc}\toprule
Phase & Nucleus & Baseline & Redesign & Speedup\\\midrule
\texttt{grstep} & $^{208}$Pb & 0.263 s @ 48 & 0.152 s @ 128 & $1.74\times$\\
 & $^{102}$Zr & 0.366 s @ 48 & 0.226 s @ 128 & $1.62\times$\\
\texttt{diagstep} & $^{208}$Pb & 0.748 s @ 192 & 0.107 s @ 192 & $6.98\times$\\
 & $^{102}$Zr & 1.260 s @ 192 & 0.188 s @ 192 & $6.71\times$\\
\texttt{pairing} & $^{102}$Zr & 0.841 s @ 48 & 0.602 s @ 64 & $1.40\times$\\
\texttt{add\_density} & $^{208}$Pb & 0.086 s @ 16 & 0.043 s @ 48 & $2.01\times$\\
 & $^{102}$Zr & 0.100 s @ 32 & 0.057 s @ 48 & $1.76\times$\\\bottomrule
\end{tabular}
\end{table}

\begin{table}[tbp]\centering
\caption{Scaling range and behaviour beyond it ($^{208}$Pb except where noted). ``Penalty'' is the time at 192 cores relative to the best time; the end-to-end row is read from Fig.~\ref{fig:e2e}.}
\label{tab:scaling}
\begin{tabular}{lcccc}\toprule
Phase & Baseline best @ & Baseline penalty & Redesign best @ & Redesign penalty\\\midrule
\texttt{grstep} & 48 & $2.4\times$ & 128 & ${\approx}1.0$\\
\texttt{diagstep} & 8 (plateau) & --- & 192 & ---\\
\texttt{pairing} ($^{102}$Zr) & 48 & $1.5\times$ & 64 & $1.2\times$\\
\texttt{add\_density} & 16 & $4.7\times$ & 48 & ${\approx}1.0$\\
End-to-end & 32 & $1.4\times$ & 128 & ${\approx}1.0$\\\bottomrule
\end{tabular}
\end{table}

\begin{table}[tbp]\centering
\caption{Iteration budget (s) at a single core count. Baseline and redesign each at their own best configuration; ``other'' is mean-field construction, single-particle properties and bookkeeping.}
\label{tab:budget}
\begin{tabular}{lcccc}\toprule
 & \multicolumn{2}{c}{$^{208}$Pb} & \multicolumn{2}{c}{$^{102}$Zr}\\\cmidrule(lr){2-3}\cmidrule(lr){4-5}
Phase & base @32 & new @128 & base @48 & new @96\\\midrule
\texttt{diagstep} & 0.83 & 0.11 & 1.39 & 0.24\\
\texttt{grstep} & 0.27 & 0.15 & 0.37 & 0.23\\
\texttt{pairing} & --- & --- & 0.84 & 0.60\\
\texttt{add\_density} & 0.09 & 0.05 & 0.12 & 0.06\\
other & 0.31 & 0.14 & 0.33 & 0.16\\\midrule
total & 1.50 & 0.45 & 3.05 & 1.28\\\bottomrule
\end{tabular}
\end{table}

\subsection{The capacity relation: footprint and knee}\label{sec:capacity}
Section~\ref{sec:mapping} argued that the mean-field application must be organized around the cache domain: the orbitals resident in a domain must fit its cache $L$, and this sets the team width through $w = D/n$ with occupancy $n = L/F_{\text{hot}}$ --- the domain capacity divided by one orbital's hot footprint. Two quantities were left to measure here: the hot footprint itself, and whether the penalty really begins at the domain boundary.

The hot footprint follows from a static accounting of one orbital's pass through the mean-field application, reading the array touches step by step (Fig.~\ref{fig:footprint}). An orbital's full working set is about 31~MiB, but at any single step only its live arrays are resident --- a peak of about 17.7~MiB during the derivative transforms, the rest cold streaming scratch written once. Two orbitals therefore hold roughly $2\times 17.7 \approx 35$~MiB of hot data, essentially the 32 MB domain, so a domain accommodates two orbitals and no more: $n=2$, and $w = D/n = 4$. This is precisely why $w=4$ works --- two orbitals' full \emph{allocations} far exceed a domain, but only their ${\sim}35$~MiB of hot data must stay resident.

\begin{figure}[tbp]\centering
  \includegraphics[width=\textwidth]{figure8}
  \caption{Working set of one orbital through the mean-field application at $48^3$, per step: the touched (hot) fields---state and scratch, and the shared mean-field arrays---against the cold remainder. The full working set ${\approx}$ one 32~MiB domain; the instantaneous peak is ${\sim}18$~MiB.}
  \label{fig:footprint}
\end{figure}

To confirm the boundary is the domain and not the node, we run the mean-field kernel on eight threads under two placements: \emph{spread}, one thread per domain, so each orbital has a domain to itself, and \emph{close}, all eight in one domain, so eight orbitals share it. Sweeping the grid scales the per-orbital hot footprint (Table~\ref{tab:knee}), and the close/spread time ratio is the knee (Fig.~\ref{fig:cap}): near one while eight orbitals still fit the domain, it rises sharply once they do not. With eight orbitals in a 32 MB domain the crossing falls at ${\sim}4$~MiB per orbital --- between $24^3$ (2.2~MiB, ratio 1.06) and $32^3$ (5.2~MiB, ratio 2.95) --- after which the ratio climbs to 4.2 and saturates near 4.4.

\begin{figure}[tbp]\centering
  \includegraphics[width=\textwidth]{figure7}
  \caption{Direct test of the capacity relation (Section~\ref{sec:capacity}). Left: mean-field time versus per-domain occupancy at fixed grid, a $4.2\times$ rise from packing alone at fixed cores, work and aggregate cache. Right: close/spread time ratio as the grid, and thus the co-resident footprint, is swept---flat while the footprint fits the 32~MiB domain, rising once it exceeds it, steepest as it sweeps through capacity.}
  \label{fig:cap}
\end{figure}

The penalty is thus gated by per-domain occupancy, not by the node's aggregate cache, which at 768 MB is ample. The binding capacity is $L$ per domain, exactly as Eq.~\eqref{eq:cap} assumes, and the knee falls at the domain size the static footprint predicts --- closing the loop between the ${\sim}18$~MiB hot footprint, two orbitals per 32 MB domain, and the team width $w=4$.

\begin{table}[tbp]\centering
\caption{The knee. Mean-field time under spread vs.\ close placement of eight threads as the grid is swept; the per-orbital hot footprint scales with the grid (Fig.~\ref{fig:footprint}), and the close/spread ratio departs from one once eight orbitals' hot data exceed the 32 MB domain.}
\label{tab:knee}
\begin{tabular}{lcccc}\toprule
grid & hot / orbital (MiB) & spread (s) & close (s) & close/spread\\\midrule
$24^3$ & 2.2 & 0.081 & 0.087 & 1.06\\
$32^3$ & 5.2 & 0.154 & 0.455 & 2.95\\
$40^3$ & 10.3 & 0.372 & 1.382 & 3.72\\
$48^3$ & 17.7 & 0.742 & 3.123 & 4.21\\
$56^3$ & 28.1 & 1.301 & 5.756 & 4.43\\\bottomrule
\end{tabular}
\end{table}

\subsection{Portability across cache organizations}\label{sec:port}
Figure~\ref{fig:port} repeats the end-to-end comparison on the two secondary machines of Table~\ref{tab:machines} --- a 128-core Zen 2 node and a 40-core Skylake node. This is not a ranking of the machines: they differ in core count and memory system, so their absolute speeds say nothing about the algorithms. What matters is that the three platforms differ in precisely the property Eq.~\eqref{eq:width} singles out --- Skylake has no small private domain (its last-level cache is shared across a socket)~\cite{intel_skylake}, Zen 2 has the smallest and most expensively bridged domains~\cite{amd_rome}, and Zen 4 is intermediate --- and on all three the baseline shows the characteristic upturn while the redesign does not.

We used the same team width $w=4$ on all three machines, without re-tuning. That it holds across cache organizations as different as a socket-wide shared L3 and a four-core private domain is a robustness result: the width need not be re-tuned per machine. It is also consistent with Eq.~\eqref{eq:width}, whose bound simply stops binding when the last-level cache is shared across a socket ($L$ is large), leaving $w$ to be set by transform-batch efficiency instead. A width sweep on the secondary machines would test this more sharply and is left to future work.

\begin{figure}[tbp]\centering
  \includegraphics[width=\textwidth]{figure9}
  \caption{The same behaviour on two further platforms: a 128-core Zen~2 node (left) and a 40-core Skylake node (right), dashed for the baseline and solid for the redesign. The same team width $w=4$ is used on all three machines without retuning.}
  \label{fig:port}
\end{figure}

\subsection{Correctness and reproducibility}\label{sec:correct}
The redesigned solver reproduces the reference converged energies to the printed precision --- $-1635.68$~MeV for $^{208}$Pb and $-859.68$~MeV for $^{102}$Zr --- with matching fluctuations along the iteration history. The convergence measure $\Delta S$ of Eq.~\eqref{eq:ds} behaves identically, an independent check that the reformulated constraint step is numerically equivalent to the original and not merely close.

Results are bit-identical at fixed thread count and team width. They are not bit-identical across different thread counts, because the order of the block reductions changes; this is the usual behaviour of blocked floating-point reductions, and the variation is at the level of rounding. The output-owned density merge of Section~\ref{sec:dens} is what makes even the fixed-configuration guarantee attainable, since it removes the runtime's tree reduction and its implementation-defined ordering.

% ======================================================================
\section{Conclusions and outlook}\label{sec:concl}
The concurrency HFBFFT inherited comes from one effective axis, the orbital index --- fixed by the nucleus, absent from the two dominant phases, and unable to grow with the machine. This paper's reformulations shift the solver's parallelism onto axes that instead grow with the numerical resolution, and that shift is what recovers scaling: the solver keeps improving across a full 192-core node and over the socket boundary, where the baseline saturates at 32--48 cores, cutting per-iteration time by roughly $3.4\times$ for $^{208}$Pb and $2.4\times$ for $^{102}$Zr (Table~\ref{tab:budget}) with converged energies unchanged and the same configuration holding on three different cache organizations.

Several directions remain. The present work is shared-memory only, and the axes exposed here are exactly those a distributed-memory layer could use; combining cooperative teams with distribution over orbitals is the natural next step, though it must contend with the state-dependent Hamiltonian of Eq.~\eqref{eq:mf}, which needs the full matrix $\lambda$ --- a difficulty the HF+BCS formulation of Sky3D avoids. The pairing sub-iteration remains the ceiling for open-shell nuclei, where only a reformulation that exposes an axis in the recurrence would help. GPU offload suits the slab decomposition and the blocked kernels, though not the recurrence; and the team width could be chosen automatically from runtime topology through Eq.~\eqref{eq:width}.

\section*{Acknowledgements}
[To be completed.]

% ======================================================================
\begin{thebibliography}{99}
\bibitem{hfbfft} M.~Chen, T.~Li, B.~Schuetrumpf, P.-G.~Reinhard, W.~Nazarewicz, Three-dimensional Skyrme Hartree--Fock--Bogoliubov solver in coordinate-space representation, Comput. Phys. Commun. 276 (2022) 108344.
\bibitem{sky3d} J.~A.~Maruhn, P.-G.~Reinhard, P.~D.~Stevenson, A.~S.~Umar, The TDHF code Sky3D, Comput. Phys. Commun. 185 (2014) 2195--2216.
\bibitem{sky3dmpi} M.~Afibuzzaman, B.~Schuetrumpf, H.~M.~Aktulga, Scalable nuclear density functional theory with Sky3D, Comput. Phys. Commun. 223 (2018) 34--44.
\bibitem{hfbtho} M.~V.~Stoitsov, N.~Schunck, M.~Kortelainen, N.~Michel, H.~Nam, E.~Olsen, J.~Sarich, S.~Wild, Axially deformed solution of the Skyrme--Hartree--Fock--Bogoliubov equations using the transformed harmonic oscillator basis (II) HFBTHO v2.00d, Comput. Phys. Commun. 184 (2013) 1592--1604.
\bibitem{hfodd} N.~Schunck, J.~Dobaczewski, J.~McDonnell, W.~Satu{\l}a, J.~A.~Sheikh, A.~Staszczak, M.~Stoitsov, P.~Toivanen, Solution of the Skyrme--Hartree--Fock--Bogolyubov equations in the Cartesian deformed harmonic-oscillator basis (VII) HFODD (v2.49t), Comput. Phys. Commun. 183 (2012) 166--192.
\bibitem{bender} M.~Bender, P.-H.~Heenen, P.-G.~Reinhard, Self-consistent mean-field models for nuclear structure, Rev. Mod. Phys. 75 (2003) 121--180.
\bibitem{blum} V.~Blum, G.~Lauritsch, J.~A.~Maruhn, P.-G.~Reinhard, Comparison of coordinate-space techniques in nuclear mean-field calculations, J. Comput. Phys. 100 (1992) 364--376.
\bibitem{fftw} M.~Frigo, S.~G.~Johnson, FFTW: an adaptive software architecture for the FFT, Proc. IEEE Int. Conf. Acoust. Speech Signal Process. 3 (1998) 1381--1384.
\bibitem{vautherin1972} D.~Vautherin, D.~M.~Brink, Hartree--Fock calculations with Skyrme's interaction. I. Spherical nuclei, Phys. Rev. C 5 (1972) 626--647.
\bibitem{chabanat1998} E.~Chabanat, P.~Bonche, P.~Haensel, J.~Meyer, R.~Schaeffer, A Skyrme parametrization from subnuclear to neutron star densities. Part II. Nuclei far from stabilities, Nucl. Phys. A 635 (1998) 231--256.
\bibitem{ringschuck} P.~Ring, P.~Schuck, The Nuclear Many-Body Problem, Springer-Verlag, Berlin, 1980.
\bibitem{dobaczewski1984} J.~Dobaczewski, H.~Flocard, J.~Treiner, Hartree--Fock--Bogolyubov description of nuclei near the neutron-drip line, Nucl. Phys. A 422 (1984) 103--139.
\bibitem{erler2012} J.~Erler, N.~Birge, M.~Kortelainen, W.~Nazarewicz, E.~Olsen, A.~M.~Perhac, M.~Stoitsov, The limits of the nuclear landscape, Nature 486 (2012) 509--512.
\bibitem{ballard2011} G.~Ballard, J.~Demmel, O.~Holtz, O.~Schwartz, Minimizing communication in numerical linear algebra, SIAM J. Matrix Anal. Appl. 32 (2011) 866--901.
\bibitem{amdgenoa} Advanced Micro Devices, AMD EPYC 9004 Series Processors (``Genoa'', Zen 4) product specifications, 2022.
\bibitem{openblas} Z.~Xianyi, W.~Qian, Z.~Yunquan, Model-driven Level-3 BLAS performance optimization, in: IEEE ICPADS, 2012; OpenBLAS, \url{https://www.openblas.net}.
\bibitem{intel_skylake} A.~Kumar, Intel Xeon Processor Scalable Family (Skylake-SP) architecture, in: IEEE Hot Chips 29 Symposium (HCS), 2017.
\bibitem{intel_haswell} Intel Corporation, Intel Xeon Processor E5-2698 v3 product specifications, 2014.
\bibitem{amd_rome} Advanced Micro Devices, AMD EPYC 7002 Series Processors (``Rome'', Zen 2), 2019.
\bibitem{blis-mt} T.~M.~Smith, R.~van de Geijn, M.~Smelyanskiy, J.~R.~Hammond, F.~G.~Van Zee, Anatomy of high-performance many-threaded matrix multiplication, in: Proc. 2014 IEEE 28th Int. Parallel and Distributed Processing Symposium (IPDPS), 2014, pp. 1049--1059.
\end{thebibliography}

\end{document}
